% Part: propositional-logic
% Chapter: syntax-and-semantics
% Section: formulas

\documentclass[../../../include/open-logic-section]{subfiles}

\begin{document}

\olfileid{pl}{syn}{fml}

\olsection{بياني \usetoken{P}{formula}}

د بياني منطق !!^{formula}s له \emph{!!{propositional
    variable}s}\iftag{prvFalse}{\iftag{prvTrue}{،}{ او} د دروغوالي له
  بياني ثابت~$\lfalse$}{}\iftag{prvTrue}{\iftag{prvFalse}{
    او}{} د رښتياوالي له بياني ثابت~$\ltrue$}{} څخه د \emph{منطقي
  نښلوونکو} په کارولو جوړېږي۔

\begin{enumerate}
\item د !!{propositional variable}s $\Obj p_0$،
  $\Obj p_1$، \dots يو !!^a{denumerable} سټ~$\PVar$
\tagitem{prvFalse}{د !!{falsity} بياني ثابت~$\lfalse$۔}{}
\tagitem{prvTrue}{د !!{truth} بياني ثابت~$\ltrue$۔}{}
\item منطقي نښلوونکي:
  \startycommalist
  \iftag{prvNot}{\ycomma $\lnot$ (نفي)}{}%
  \iftag{prvAnd}{\ycomma $\land$ (عطف)}{}%
  \iftag{prvOr}{\ycomma $\lor$ (فصل)}{}%
  \iftag{prvIf}{\ycomma $\lif$ (!!{conditional})}{}%
  \iftag{prvIff}{\ycomma $\liff$ (!!{biconditional})}{}%
\item د وقف نښې: (، ) او کامه۔
\end{enumerate}

د بياني منطق دا ژبه په $\Lang L_0$ ښيو۔

\iftag{defNot,defOr,defAnd,defIf,defIff,defTrue,defFalse,defEx,defAll}{%
د پورته معرفي شويو بنسټيزو نښلوونکو تر څنګ، دا لاندې
\emph{تعريف شوې} نښې هم کاروو:
\startycommalist
  \iftag{defNot}{\ycomma $\lnot$ (نفي)}{}%
  \iftag{defAnd}{\ycomma $\land$ (عطف)}{}%
  \iftag{defOr}{\ycomma $\lor$ (فصل)}{}%
  \iftag{defIf}{\ycomma $\lif$ (!!{conditional})}{}%
  \iftag{defIff}{\ycomma $\liff$ (!!{biconditional})}{}%
  \iftag{defFalse}{\ycomma $\lfalse$ (!!{falsity})}{}%
  \iftag{defTrue}{\ycomma $\ltrue$ (!!{truth})}}{}۔

\begin{tagblock}{defNot,defOr,defAnd,defIf,defIff,defTrue,defFalse,defEx,defAll}
\begin{explain}
تعريف شوې نښه په رسمي توګه د ژبې برخه نۀ وي، خو د يوې غير رسمي
لنډې نښې په توګه معرفي کېږي: په دې سره هغه فارمولونه لنډولے شو چې
که يوازې بنسټيزې نښې مو کارولې، ډېر اوږد به شوي وو۔ دا ئې څرګنده
ګټه ده۔ خو لويه ګټه ئې دا ده چې ثبوتونه لنډېږي۔ که کومه نښه
بنسټيزه وي، بايد په ثبوتونو کښې جلا وڅېړل شي۔ نو څومره چې بنسټيزې
نښې ډېرې وي، هومره به زموږ ثبوتونه اوږده وي۔
\end{explain}
\end{tagblock}

% Alternate symbols

\begin{intro}
کېدے شي تاسو له هغو اصطلاحاتو او نښو سره بلد يئ چې زموږ له پورته
کارولو سره توپير لري۔ د منطق متنونه، او ښوونکي، عموماً د ``نفي''
دپاره $\sim$، $\neg$ او~!\، او د ``عطف'' دپاره $\wedge$، $\cdot$ او
$\&$ کاروي۔ د ``شرطي'' يا ``استلزام'' عامې نښې $\rightarrow$،
$\Rightarrow$ او $\supset$ دي۔
\iftag{prvIff,defIff}{د ``دوه اړخيز شرطي''، ``دوه اړخيز استلزام''
  يا ``(مادي) هم ارزښت'' نښې $\leftrightarrow$،
  $\Leftrightarrow$ او $\equiv$ دي۔}{}
\iftag{prvFalse,defFalse}{د $\lfalse$ نښې ته په بېلو ځايونو کښې
  ``دروغوالے''، ``فالسوم''، ``محال'' يا ``ښکته'' ويل کېږي۔}{}
\iftag{prvTrue,defTrue}{د $\ltrue$ نښې ته په بېلو ځايونو کښې
  ``رښتياوالے''، ``ويروم'' يا ``پورته'' ويل کېږي۔}{}
\end{intro}

\begin{defn}[فارمول]
\ollabel{defn:formulas}
د بياني منطق د \emph{!!{formula}s} سټ~$\Frm[L_0]$ په استقرايي ډول
داسې تعريفېږي:
\begin{enumerate}
\tagitem{prvFalse}{$\lfalse$ يو اټومي !!{formula} دے۔}{}

\tagitem{prvTrue}{$\ltrue$ يو اټومي !!{formula} دے۔}{}

\item هر !!{propositional variable}~$\Obj p_i$ يو اټومي
  !!{formula} دے۔

\tagitem{prvNot}{که $!A$ يو !!a{formula} وي، نو $\lnot !A$ هم
  يو !!a{formula} دے۔}{}

\tagitem{prvAnd}{که $!A$ او $!B$ !!{formula}s وي، نو $(!A \land
  !B)$ يو !!a{formula} دے۔}{}

\tagitem{prvOr}{که $!A$ او $!B$ !!{formula}s وي، نو $(!A \lor !B)$
  يو !!a{formula} دے۔}{}

\tagitem{prvIf}{که $!A$ او $!B$ !!{formula}s وي، نو $(!A \lif !B)$
  يو !!a{formula} دے۔}{}

\tagitem{prvIff}{که $!A$ او $!B$ !!{formula}s وي، نو $(!A \liff !B)$
  يو !!a{formula} دے۔}{}

\tagitem{limitClause}{بل هېڅ شے !!a{formula} نۀ دے۔}{}
\end{enumerate}
\end{defn}

\begin{explain}
د !!{formula}s تعريف يو \emph{استقرايي تعريف} دے۔ په اصل کښې موږ
د !!{formula}s سټ په نامتناهي ډېرو پړاوونو کښې جوړوو۔ په لومړني
پړاو کښې ټول اټومي فارمولونه فارمولونه ګڼو؛ دا د تعريف له لومړنيو
څو حالتونو سره سمون لري، يعنې د
\iftag{prvTrue}{$\ltrue$، }{}%
\iftag{prvFalse}{$\lfalse$، }{}%
$\Obj p_i$ حالتونه۔ په دې توګه ``اټومي !!{formula}'' د همدې بڼې هر
!!{formula} ته وايي۔

د تعريف نور حالتونه له هغو !!{formula}s څخه د نويو !!{formula}s د
جوړولو قاعدې راکوي چې له مخکښې جوړ شوي وي۔ په دويم پړاو کښې له دغو
قاعدو څخه په کار اخيستو له اټومي !!{formula}s څخه !!{formula}s
جوړولے شو۔ په درېيم پړاو کښې له اټومي فارمولونو او په دويم پړاو کښې
له ترلاسه شويو فارمولونو څخه نوي فارمولونه جوړوو، او همداسې مخکښې
ځو۔ !!{formula} هر هغه شے دے چې بالاخره په داسې کوم پړاو کښې جوړ
شي، او بل هېڅ شے نۀ دے۔
\end{explain}

کله چې د دوه ځايه نښلوونکي~$\ast$ په کارولو له $!B$ او $!C$ څخه
جوړ شوے فارمول $(!B \ast !C)$ ليکو، ډېر ځله د قوسونو تر ټولو بهرنی
جوړ پرېږدو او يوازې~$!B \ast !C$ ليکو۔

\begin{tagblock}{defNot,defOr,defAnd,defIf,defIff,defTrue,defFalse,defEx,defAll}
\begin{defn}
په تعريف شويو عملګرو جوړ شوي فارمولونه بايد داسې وپوهول شي:

\begin{tagenumerate}{defTrue,defFalse,defNot,defOr,defAnd,defIf,defIff,defEx,defAll}

\tagitem{defTrue}{$\ltrue$ د دې لنډه نښه ده:
  \iftag{prvFalse}{$\lnot\lfalse$}{$(!A \lor \lnot !A)$ د کوم
    ټاکلي اټومي !!{formula}~$!A$ دپاره}۔}{}

\tagitem{defFalse}{$\lfalse$ د دې لنډه نښه ده:
  \iftag{prvTrue}{$\lnot\ltrue$}{$(!A \land \lnot !A)$ د کوم
    ټاکلي اټومي !!{formula}~$!A$ دپاره}۔}{}

\tagitem{defNot}{$\lnot !A$ د $!A \lif \lfalse$ لنډه نښه ده۔}{}

\tagitem{defOr}{$!A \lor !B$ د دې لنډه نښه ده:
  \iftag{prvAnd}{$\lnot(\lnot !A \land \lnot !B)$}{$\lnot !A \lif
    !B$}۔}{}

\tagitem{defAnd}{$!A \land !B$ د دې لنډه نښه ده:
  \iftag{prvOr}{$\lnot(\lnot !A \lor \lnot !B)$}{$\lnot (!A \lif
    \lnot !B)$}۔}{}

\tagitem{defIf}{$!A \lif !B$ د دې لنډه نښه ده:
  \iftag{prvOr}{$\lnot !A \lor !B$}{$\lnot (!A \land \lnot !B)$}۔}{}
\textbf{د ژباړې د نسخې سمون (OLPL-001):} د سرچينې د لومړي بديل
په ښۍ خوا کښې يو تړلے قوس بې له پرانيستي قوسه راغلے دے۔ دلته هغه
زياتي قوس لرې شوے، څو لنډه نښه هم نحوي سمه وي او هم د فصل له
معمولي تعريف سره برابره وي۔

\tagitem{defIff}{$!A \liff !B$ د $(!A \lif !B) \land (!B
  \lif !A)$ لنډه نښه ده۔}{}
\end{tagenumerate}
\end{defn}
\end{tagblock}

\begin{defn}[نحوي يوشانوالی]
د $\ident$ نښه د نښو د توريزو لړيو ترمنځ نحوي يوشانوالی څرګندوي؛
يعنې $!A \ident !B$ هله او يوازې هله چې $!A$ او $!B$ د نښو د يو
شان اوږدوالي توريزې لړۍ وي او په هر ځاے کښې يوه نښه ولري۔
\end{defn}

د $\ident$ نښې دواړو خواوو ته د يو ځاے کولو له لارې ترلاسه شوې
توريزې لړۍ هم راتلے شي۔ مثلاً $!A \ident (!B \lor !C)$ دا معنٰی لري:
د نښو توريزه لړۍ~$!A$ هماغه توريزه لړۍ ده چې په همدې ترتيب د پرانيستي
قوس، د $!B$ توريزې لړۍ، د $\lor$ نښې، د~$!C$ توريزې لړۍ او د تړلي
قوس په يو ځاے کولو ترلاسه کېږي۔ که داسې وي، نو پوهېږو چې د $!A$
لومړۍ نښه پرانيستے قوس دے، $!A$ د $!B$ توريزه لړۍ د فرعي توريزې
لړۍ په توګه لري چې له دويمې نښې پيلېږي، ورپسې $\lor$ راځي، او داسې نور۔

\end{document}
